
\section{Chaining theoretical analysis}
\label{sec:chaining-analysis-two}

We outlined the minimap2 algorithms for chaining, defined scores formally, but I would like to spend additional time on more formal analysis of chain properties.
There is extensive literature on analysis of theoritcal and practical aligners performance, however I wanted to focus on exact minimap2 scoring and chaining model. It also serves as an additional resource for understanding pipeline steps presented so far.

\paragraph{Point substitutions.}
Set \(h=k+w-1\). A span of \(h\) bases contains one complete window of \(w\)
candidate \(k\)-mers. For this bound, assume minimap2 uses
non-homopolymer-compressed (non-HPC) minimizers and odd \(k\). Let's suppose aligned
bases are changed independently with some 
probability \(\varepsilon\), with no insertions or deletions. A changed base is
replaced by one of the other three bases. Assume the fragment lies in an
ambiguity-free record, so a stored window minimum eventually leaves the
rolling window or is emitted at the end of the record. An unchanged span of
\(h\) bases gives both sequences the same complete minimizer window. Let
\(A_L^{\mathrm{raw}}\) count shared minimizers before occurrence filtering and
masking of low-complexity sequence. Partitioning a homologous fragment of length \(L\)
into disjoint spans gives
\begin{equation}
  \Pr\!\left(A_L^{\mathrm{raw}}=0\right)
  \leq
  \left[1-(1-\varepsilon)^h\right]^{\lfloor L/h\rfloor}.
  \label{eq:chaining-no-anchor-substitution-bound}
\end{equation}
This bound does not treat overlapping minimizers as independent. The
probability of observing no raw anchor decreases exponentially with
\(\lfloor L/h\rfloor\) under this substitution model. Extending the bound to
retained anchors requires a condition that ensures the shared minimizer passes
the occurrence filters.

\paragraph{Average-case anchor population.}
Equation~\ref{eq:chaining-no-anchor-substitution-bound} concerns one fixed
fragment. Shaw and Yu study the number of anchors in a random reference-query
pair~\cite{shaw_yu_seed_chain_extend_2023}. Their reference contains
\(n_R\) possible \(k\)-mer starts and is sampled uniformly and independently
over the DNA alphabet. The query contains \(n_Q\leq n_R\) possible starts and
is copied from a homologous reference substring. Each query base is then
substituted independently with probability \(0\leq\varepsilon<1\). A
substitution changes the base to one of the other three bases. The model has
no insertions or deletions.

Let \(N_{\mathrm{hom}}\) count exact \(k\)-mer matches at homologous starts and
let \(N_{\mathrm{off}}\) count matches between all other pairs of starts. Set
\begin{equation}
  \begin{aligned}
    q_k&=(1-\varepsilon)^k,
    &s_k&=4^{-k},\\
    \mathbb E[N_{\mathrm{hom}}]&=n_Qq_k,
    &\mathbb E[N_{\mathrm{off}}]&=n_Q(n_R-1)s_k.
  \end{aligned}
  \label{eq:chaining-average-anchor-counts}
\end{equation}
The second line is Theorem~4 of Shaw and Yu for one forward
orientation~\cite{shaw_yu_seed_chain_extend_2023}.
The first term counts surviving homologous \(k\)-mers. The second term counts
accidental exact matches away from the homologous coordinate. These
expectations count all exact \(k\)-mer-match pairs before minimizer selection,
masking, and occurrence filtering.

Overlapping \(k\)-mers make neighbouring anchor indicators dependent. Set
\(N=N_{\mathrm{hom}}+N_{\mathrm{off}}\) and
\(\bar\mu=n_Qq_k+n_Qn_Rs_k\). Then
\(\mathbb E[N]\leq\bar\mu\). Combining the three bounds in Lemma~3 and using
\(n_R-1\leq n_R\) gives~\cite{shaw_yu_seed_chain_extend_2023}
\begin{equation}
  \mathbb E[N^2]\leq B_k,
  \qquad
  B_k=\bar\mu^2+2kn_Qq_k+8k(k+1)n_Qn_Rs_k.
  \label{eq:chaining-average-anchor-second-moment}
\end{equation}
This bound retains the dependence between overlapping anchors.

The same model also bounds long intervals without a homologous exact
\(k\)-mer anchor. Set \(k=C\log_4n_R\), with \(C>0\), for input sizes where
\(k\) is an integer, and let \(\xi=-\log_4(1-\varepsilon)\). If
\(G_{\max}\) is the longest run of homologous \(k\)-mer starts without an
anchor, Lemma~6 gives~\cite{shaw_yu_seed_chain_extend_2023}
\begin{equation}
  \begin{aligned}
    g(n_R)&=\frac{50C}{8}\,
      \log_4 n_R\,\ln n_R\,n_R^{C\xi},\\
    \Pr\!\left(G_{\max}\leq g(n_R)\right)&\geq1-\frac{1}{n_R}.
  \end{aligned}
  \label{eq:chaining-average-largest-gap}
\end{equation}
The corresponding base interval is at most \(g(n_R)+k-1\) bases long.

For the chaining and asymptotic-work conclusions below, impose the scaling
assumptions \(0\leq\xi<1\),
\(n_Q=\Omega(n_R^{2C\xi+\tau})\) for some \(\tau>0\), and
\(C>1/(1-\xi)\). These assumptions require \(2C\xi+\tau\leq1\) to remain
compatible with \(n_Q\leq n_R\).

\paragraph{Consequences for minimap2 work.}
Let \(M\) be the total number of retained anchors passed to one minimap2
non-RMQ \texttt{mg\_lchain\_dp} call, including both canonical orientation
channels. Write
\(\mu_{\mathrm{mm}}=\mathbb E[M]\),
\(\nu_{\mathrm{mm}}=\mathbb E[M^2]\), and
\(I=\texttt{max\_iter}\geq0\). Inspection of the minimap2 loops and then
taking expectations gives the following bounds for
Algorithms~\ref{alg:chaining-b1} and~\ref{alg:chaining-minimap2-backtrack}:
\begin{align}
  W_{\mathrm{B1}}
    &=O\!\left(M\bigl(1+\min\{I,M\}\bigr)\right),
  &W_{\mathrm{bt}}&=O(M^2),\notag\\[-2pt]
  \mathbb E[W_{\mathrm{B1}}]
    &=O\!\left((I+1)\mu_{\mathrm{mm}}\right),
  &\mathbb E[W_{\mathrm{bt}}]&=O(\nu_{\mathrm{mm}}).
  \label{eq:chaining-minimap2-work-by-anchor-count}
\end{align}
For fixed \(I\), expected forward work is \(O(\mu_{\mathrm{mm}})\), and
expected scalar backtracking is \(O(\nu_{\mathrm{mm}})\). No independence
assumption is needed for either statement.

The Shaw-Yu moment bounds apply to one part of \(M\). Let \(M_+\) count
retained non-HPC minimizer anchors in the forward channel, where the reference
and query canonical strand bits agree. With collision-free hashing, equal keys
in this channel imply equal original \(k\)-mers, so \(M_+\leq N\) on every
input. Minimizer selection, masking, and the occurrence threshold used by
either collection pass can only remove pairs from the exact-match set counted
by \(N\). Equations~\ref{eq:chaining-average-anchor-counts}
and~\ref{eq:chaining-average-anchor-second-moment} therefore give
\begin{equation}
  \mathbb E[M_+]\leq\bar\mu,
  \qquad
  \mathbb E[M_+^2]\leq B_k.
  \label{eq:chaining-average-minimap2-oriented-moments}
\end{equation}
Under these scaling assumptions,
\(\mathbb E[M_+]=O(n_Qn_R^{-C\xi})\) and
\(\mathbb E[M_+^2]=O(n_Q^2n_R^{-2C\xi})\).
Minimap2 also passes opposite-orientation anchors to
\texttt{mg\_lchain\_dp}. Shaw and Yu do not bound the first and second moments
of that channel or its cross moment with \(M_+\). Thus, the \(M_+\) bounds in
Equation~\ref{eq:chaining-average-minimap2-oriented-moments} do not bound
\(\mu_{\mathrm{mm}}\) or \(\nu_{\mathrm{mm}}\).

For comparison, Shaw and Yu analyse an exact optimal linear-gap chainer with
per-instance time \(O(N\log N)\). Under those scaling assumptions, their
Theorem~6~\cite{shaw_yu_seed_chain_extend_2023} proves
\(\mathbb E[N\log N]=O(n_Qn_R^{-C\xi}\log n_Q)\). This
upper-bounds the expected time of their chainer. It does not give the expected
running time of minimap2's B1 score or its backtracking procedure.

The transfer to \(M_+\) uses only Shaw and Yu's all-\(k\)-mer moment bounds.
Their sketch-specific proof does not cover
minimizers~\cite{shaw_yu_seed_chain_extend_2023}, and their optimal linear-gap
chainer differs from minimap2's score and bounded predecessor search. The
random model omits indels and structured genomic repeats. It also omits
occurrence caps and the fixed \(k\) used by minimap2 in practice. The
dense-filler construction below shows that scalar backtracking can still take
time quadratic in the number of retained anchors.

\paragraph{Dense fillers with no shared values.}
Let \(U\) and \(V\) be independent uniform DNA fillers. They contain
\(N_U=|U|-k+1\) and \(N_V=|V|-k+1\) candidate \(k\)-mers. For two independent
canonical \(k\)-mers, the equality probability is
\begin{equation}
  p_k=
  \begin{cases}
    2\,4^{-k}, & k\text{ odd},\\[2pt]
    2\,4^{-k}-4^{-3k/2}, & k\text{ even}.
  \end{cases}
  \label{eq:chaining-canonical-collision-probability}
\end{equation}
Let \(A^{\mathrm{raw}}(U,V)\) be the number of raw anchor pairs between the two
fillers. For \(k=15\), minimap2's masked \texttt{hash64} is injective on the
30-bit canonical encodings. Every raw minimizer anchor must therefore come
from equal canonical \(k\)-mers. A union bound over all candidate pairs gives
\begin{equation}
  \Pr\!\left(A^{\mathrm{raw}}(U,V)>0\right)\leq N_UN_Vp_k.
  \label{eq:chaining-unrelated-filler-bound}
\end{equation}
If \(N_UN_Vp_k\leq\eta\), the fillers have no shared value with probability at
least \(1-\eta\). For the default non-HPC odd \(k=15\), start with a clean
rolling buffer and follow the filler with \(w\) ambiguous bases. After
excluding the final \(w\) candidate starts, filler \(U\) emits at least
\(\max\{0,\lfloor(N_U-w)/w\rfloor\}\) selected positions. The same bound
holds for \(V\).
Equation~\ref{eq:chaining-unrelated-filler-bound} uses every candidate pair, so
it does not require independent minimizer choices. Occurrence filtering can
only remove anchors from this count.

Put at least \(w\) ambiguous bases between consecutive fragments. Minimap2 restarts
\(k\)-mer formation at the boundary, and the \(w\) invalid entries also clear
its rolling minimizer buffer. No \(k\)-mer crosses the boundary. The restart
affects \(O(k+w)\) bases, which is constant for fixed \(k\) and \(w\).

\paragraph{One emitted value per conserved block.}
To discard the final minimum without emitting it, append at least \(w\) ambiguous bases
to each conserved block. The following construction applies to the default
non-HPC parameters \(k=15\) and \(w=10\). Let \(s\) be an 11-base word and let
\(B(s)\) be the first 25 bases of the periodic sequence \(ssss\cdots\). Its
eleven candidate 15-mers are \(q_0,\ldots,q_{10}\).

Choose \(G\) independent uniform 11-base words. Within one block, two
different forward candidates are equal with probability \(4^{-10}\). A
candidate cannot equal the reverse complement of another candidate. On an
odd cycle, such an alignment would pair one base with its own complement.
Across two independently chosen blocks, a pair of canonical candidates is
equal with probability at most \(2\,4^{-11}\). A union bound gives
\begin{equation}
  \Pr(\mathcal E_D)\geq
  1-G\binom{11}{2}4^{-10}
  -G(G-1)11^2 4^{-11},
  \label{eq:chaining-conserved-block-existence}
\end{equation}
where \(\mathcal E_D\) says that all \(11G\) canonical candidates are
different. The right side is positive for \(G\leq185\). A suitable fixed set
of blocks therefore exists.

For each chosen block, rotate \(s\) until its smallest hash is \(q_{10}\).
Let \(u\) be the smallest hash among \(q_0,\ldots,q_9\), and let
\(v=q_{10}\). All hashes are different under \(\mathcal E_D\), so \(v<u\).
The first complete window has the unique minimum \(u\). Candidate \(v\) then
causes \texttt{mm\_sketch} to emit \(u\) and to store \(v\) as the new
minimum. The next ten ambiguous bases overwrite the buffer after
the valid-run length has been reset to zero. They discard \(v\) without
emitting it, so the block emits
exactly \(u\). Append the separator to the final block as well, which prevents
the end-of-sequence step from emitting \(v\).

Corresponding reference and query copies of a block emit the same value with
the same canonical orientation. Values from different blocks are distinct.
Each value occurs once in the target and passes an occurrence filter that
retains singleton values, including the default \texttt{map-ont} filter. This
assumes that the constructed target is the complete indexed target. With a
background reference, the same values must also be absent from that reference.

\paragraph{Long padding with no anchors.}
The score-reducing path needs long coordinate gaps. If every gap base were
sampled independently, the union-bound error term would grow with total gap
length. Use periodic padding with \(r=k+1\). Choose one target period and one
query period of length \(r\). Repeating a period to any length uses at most \(r\) canonical
\(k\)-mer values. The periods can be chosen so that their value sets are
disjoint from each other and from the conserved values in the opposite
sequence. For independently sampled periods, a union bound gives
\begin{equation}
  \Pr(\mathcal E_P)\geq
  1-\bar p_k\left(r^2+2rG\right),
  \label{eq:chaining-periodic-padding-existence}
\end{equation}
Here, \(\mathcal E_P\) is the event that the target and query period value
sets are disjoint from each other and from conserved values in the opposite
sequence. Also, \(\bar p_k=2\,4^{-k}\). This bound is positive for \(k=15\),
\(r=16\), and \(G\leq185\), so a fixed pair of periods exists. Long fragments
made from these periods contain selected minimizers but create no target-query
anchors. They provide the large coordinate gaps in the prefix and the
score-reducing path.
Periodic padding serves only to construct a finite worst-case instance. It
does not model an independent random genome or the probability of observing
this construction in biological data.

After fixing the conserved blocks and periodic padding, sample every random
filler independently. Let
\(N_T^\circ,N_Q^\circ\) count candidate starts wholly inside the random
target and query fillers. Let \(D_T,D_Q\) count the distinct fixed values
that can be emitted by the conserved blocks and periodic padding. A union
bound over every
target-query candidate pair in the chaining group gives
\begin{equation}
\begin{aligned}
  \Pr(\mathcal E_F)\geq 1
  &-p_kN_T^\circ N_Q^\circ\\
  &-\bar p_k\!\left(N_T^\circ D_Q+D_TN_Q^\circ\right),
\end{aligned}
\label{eq:chaining-complete-filler-event}
\end{equation}
where \(\mathcal E_F\) means that no random filler value creates an unintended
anchor. The distinctness conditions exclude every fixed-value anchor except
the constructed anchor pairs. This model covers all candidate pairs in the chaining group, including pairs that
would later fail a geometric test.
Choose fixed blocks and periods that satisfy \(\mathcal E_D\) and
\(\mathcal E_P\) before sampling the fillers. If the right side of
Equation~\ref{eq:chaining-complete-filler-event} is positive, at least one
complete filler realization has only the intended anchors.

For a concrete finite check, take \(\ell=3\) and three backbone groups. Use
four random fillers per group, each with 20 candidate starts, and use the fixed
periodic padding for every longer gap. Then
\(N_T^\circ,N_Q^\circ\leq240\), \(G\leq77\), and
\(D_T,D_Q\leq93\). With \(k=15\),
Equation~\ref{eq:chaining-complete-filler-event} gives
\(\Pr(\mathcal E_F)>0.9998\).
For the same values, Equations~\ref{eq:chaining-conserved-block-existence}
and~\ref{eq:chaining-periodic-padding-existence} give lower bounds greater
than \(0.8271\) and \(0.99999\). Sampling blocks, periods, and short fillers in
that order gives a joint lower bound greater than \(0.8269\) for this finite
construction.
\coderef{minimap2/sketch.c}{Minimap2 minimizer sketch}{mm_sketch}{77}

Figure~\ref{fig:chaining-quadratic-sequence-template} shows the resulting
reference and query layout.

\begin{figure}[H]
\centering
\begin{tikzpicture}[
  x=1cm,
  y=1cm,
  font=\footnotesize,
  sequence block/.style={
    draw=black!65,
    rounded corners=1.3pt,
    minimum height=.46cm,
    inner xsep=4pt,
    inner ysep=1.5pt
  },
  sequence prefix/.style={
    sequence block,
    draw=chainGreen!80!black,
    fill=chainGreen!14
  },
  sequence backbone/.style={
    sequence block,
    draw=chainBlue!90!black,
    fill=chainBlue!15
  },
  sequence branch/.style={
    sequence block,
    draw=chainOrange!95!black,
    fill=chainOrange!17
  },
  sequence filler/.style={
    sequence block,
    draw=chainGray!85!black,
    fill=chainGray!12,
    text=black!72
  },
  sequence anchor/.style={line width=1.25pt},
  minimizer tick/.style={draw=black!48,line width=.55pt}
]
  \node[anchor=east,font=\bfseries] at (.42,4.45) {reference};
  \node[anchor=east,font=\bfseries] at (.42,2.85) {query};
  \draw[black!40] (.55,4.45)--(13.10,4.45);
  \draw[black!40] (.55,2.85)--(13.10,2.85);

  \node[sequence prefix,minimum width=.88cm] (tp) at (1.05,4.45) {$P$};
  \node[sequence prefix,minimum width=1.20cm,font=\scriptsize]
    (tg) at (2.22,4.45) {$G_t:\alpha^*$};
  \node[sequence backbone] (tcj) at (3.37,4.45) {$C_j$};
  \node[sequence filler,minimum width=.78cm] (tuj) at (4.20,4.45) {$U_j$};
  \node[sequence branch] (trj) at (5.02,4.45) {$R_j$};
  \node[sequence filler,minimum width=.72cm] (tujb) at (5.83,4.45) {$U'_j$};
  \node[sequence backbone] (tcjp) at (6.72,4.45) {$C_{j+1}$};
  \node[sequence filler,minimum width=.72cm] (tujp) at (7.64,4.45) {$U_{j+1}$};
  \node[sequence branch] (trjp) at (8.48,4.45) {$R_{j+1}$};
  \node[sequence filler,minimum width=.70cm] at (9.30,4.45) {$U'_{j+1}$};
  \node[text=black!55] at (10.07,4.45) {$\cdots$};
  \node[sequence backbone] (tcm) at (10.86,4.45) {$C_m$};
  \node[sequence filler,minimum width=.68cm] at (11.67,4.45) {$U_m$};
  \node[sequence branch] (trm) at (12.48,4.45) {$R_m$};

  \node[sequence prefix,minimum width=.88cm] (qp) at (1.05,2.85) {$P$};
  \node[sequence prefix,minimum width=1.12cm,font=\scriptsize]
    (qg) at (2.12,2.85) {$G_q:\beta^*$};
  \node[sequence backbone] (qcj) at (3.23,2.85) {$C_j$};
  \node[sequence filler,minimum width=.62cm] (qvj) at (3.93,2.85) {$V_j$};
  \node[sequence branch] (qrj) at (4.65,2.85) {$R_j$};
  \node[sequence filler,minimum width=1.02cm] (qvjb) at (5.65,2.85) {$V'_j$};
  \node[sequence backbone] (qcjp) at (6.84,2.85) {$C_{j+1}$};
  \node[sequence filler,minimum width=.86cm] (qvjp) at (7.83,2.85) {$V_{j+1}$};
  \node[sequence branch] (qrjp) at (8.87,2.85) {$R_{j+1}$};
  \node[sequence filler,minimum width=.66cm] at (9.72,2.85) {$V'_{j+1}$};
  \node[text=black!55] at (10.35,2.85) {$\cdots$};
  \node[sequence backbone] (qcm) at (11.05,2.85) {$C_m$};
  \node[sequence filler,minimum width=.62cm] at (11.82,2.85) {$V_m$};
  \node[sequence branch] (qrm) at (12.75,2.85) {$R_m$};

  \draw[sequence anchor,chainGreen!80!black] (tp.south)--(qp.north);
  \draw[sequence anchor,chainGreen!80!black] (tg.south)--(qg.north);
  \draw[sequence anchor,chainBlue!90!black] (tcj.south)--(qcj.north);
  \draw[sequence anchor,chainOrange!95!black] (trj.south)--(qrj.north);
  \draw[sequence anchor,chainBlue!90!black] (tcjp.south)--(qcjp.north);
  \draw[sequence anchor,chainOrange!95!black] (trjp.south)--(qrjp.north);
  \draw[sequence anchor,chainBlue!90!black] (tcm.south)--(qcm.north);
  \draw[sequence anchor,chainOrange!95!black] (trm.south)--(qrm.north);

  \foreach \x in {3.96,4.15,4.34,5.64,5.83,6.02,7.45,7.64,7.83,
                    9.12,9.30,9.48,11.52,11.67,11.82}
    \draw[minimizer tick] (\x,4.72)--(\x,4.90);
  \foreach \x in {3.78,3.93,4.08,5.30,5.55,5.80,6.05,7.60,7.83,8.06,
                    9.57,9.72,9.87,11.70,11.82,11.94}
    \draw[minimizer tick] (\x,3.12)--(\x,3.30);

  \draw[chainBlue!90!black,line width=1.2pt] (.72,1.92)--(1.18,1.92);
  \node[anchor=west,font=\scriptsize] at (1.27,1.92)
    {identical backbone block};
  \draw[chainOrange!95!black,line width=1.2pt] (4.75,1.92)--(5.21,1.92);
  \node[anchor=west,font=\scriptsize] at (5.30,1.92)
    {branch \(R_j\): \(\ell+1\) conserved blocks};
  \draw[minimizer tick] (9.05,1.26)--(9.05,1.64);
  \node[anchor=west,font=\scriptsize] at (9.20,1.45)
    {selected filler minimizer};
\end{tikzpicture}
\caption{Finite reference and query sequences that produce the retained
anchors used in the lower bound. \(P\) summarizes the high-scoring prefix.
\(G_t\) and \(G_q\) summarize the target and query paths that reduce the score
to \(F_C\). Their padding repeats fixed words \(\alpha\) and \(\beta\) whose
canonical \(k\)-mer sets are disjoint. Each target-query pair labeled \(C_j\)
produces backbone anchor \(c_j\). Each target-query pair labeled \(R_j\)
produces branch anchors
\(x_{j,0},\ldots,x_{j,\ell}=e_j\). Gray ticks mark filler minimizers, and
colored connectors mark retained anchors. Each small gap contains at least
\(w\) ambiguous bases. Unequal filler lengths change the diagonal between
consecutive conserved blocks.}
\label{fig:chaining-quadratic-sequence-template}
\end{figure}

Each \(C_j\) block supplies one retained anchor \(c_j\). Region \(R_j\)
contains \(\ell+1\) conserved blocks \(R_{j,0},\ldots,R_{j,\ell}\).
Equal-length target and query fillers separate consecutive blocks. These
fillers contain minimizers but produce no anchors under \(\mathcal E_F\).
The conserved blocks supply anchors
\(x_{j,0},\ldots,x_{j,\ell}=e_j\) on one diagonal. Every conserved block in
the backbone, branches, prefix, and score-reducing path uses a different
retained value under \(\mathcal E_D\).
Group \(j\) consists of backbone block \(C_j\), branch region \(R_j\), and
their adjacent fillers.

\paragraph{Symbolic score construction.}
For the \texttt{map-ont} defaults, \(k=15\), the gap coefficient is
\(0.008k=0.12\), and the distance coefficient is zero. Let
\(g(j,i)=\min\{r_i-r_j,q_i-q_j\}\). When \(g(j,i)\geq k\), a diagonal
change of \(x\) gives
\begin{equation}
  \Lambda(x)=k-
  \left\lfloor
    0.12x+\frac{1}{2}\operatorname{mglog}_{2}(x+1)
  \right\rfloor.
  \label{eq:chaining-map-ont-symbolic-link}
\end{equation}
Here, \(\operatorname{mglog}_{2}\) is minimap2's fast floating-point approximation to the
base-two logarithm. The surrounding cast truncates the result to give an integer gap penalty.
Let \(B=\texttt{bw}\), \(X=\texttt{max\_gap}\),
\(D=\texttt{max\_drop}\), \(S=\texttt{min\_sc}\), and
\(K=\texttt{min\_cnt}\). Scalar chaining sets \(D=B=500\) for these defaults.
Choose positive diagonal changes \(d_1,d_2\) and a fixed branch length
\(\ell\) such that
\begin{equation}
\begin{gathered}
  d_1+d_2\leq B<d_1+2d_2,\qquad \Lambda(d_1)=0,\\
  \Lambda(d_1+d_2)<\Lambda(d_2)<0,\\
  0<\Lambda(d_2)+\ell k<-\Lambda(d_1+d_2).
\end{gathered}
\label{eq:chaining-quadratic-link-conditions}
\end{equation}
The default \texttt{map-ont} parameter values satisfy these conditions. For
example,
\(d_1=100\), \(d_2=250\), and \(\ell=3\) give link contributions
\(0\), \(-18\), and \(-31\) for the three diagonal changes in
Equation~\ref{eq:chaining-quadratic-link-conditions}. Taking \(F_C=39\)
then gives \(F_X=21\) and \(F_E=66\), which satisfies the default
\(S=40\) and \(K=3\) thresholds.

For an anchor \(a\), let \(z(a)=r(a)-q(a)\), and let \(s_j=(-1)^j\). The
blue backbone anchors \(c_j\) and the orange branch anchors
\(x_{j,0},\ldots,x_{j,\ell}=e_j\) have
\begin{equation}
\begin{aligned}
  z(c_{j+1}) &= z(c_j)-s_jd_1,\\
  z(x_{j,t}) &= z(c_j)+s_jd_2
  && (0\leq t\leq\ell).
\end{aligned}
\label{eq:chaining-quadratic-diagonals}
\end{equation}
Consecutive anchors inside one orange branch lie on the same diagonal. One
conserved block and its separator occupy \(b_E=k+2w\) bases. A filler of
\(b_F=k+2w-1\) valid bases followed by another \(w\)-base separator contains
a complete window and emits at least one minimizer. The emitted value begins at one of positions
\(0,\ldots,w-1\) in its block. Let \(h_0\) bound coordinate offsets within a
group, let \(L\) be the spacing between backbone anchors, let \(T\) be the
branch offset from \(c_j\), and let \(\tau\) be the spacing between branch
anchors. Choose these values so that
\[
\begin{gathered}
  d_1+2d_2<h_0,\qquad h_0+4d_1<X,\\
  \tau\geq b_E+b_F+2w-1,\\
  d_2+b_E+b_F+2w-1\leq T,\\
  T+d_2+\ell\tau<h_0,\\
  \frac{X+h_0}{2}+d_1<L<X-d_1.
\end{gathered}
\]
If two emitted offsets differ, lengthen one separator by at most \(2(w-1)\)
ambiguous bases to obtain the required coordinate difference. This does not
change the random-filler counts.
Choose the coordinates of \(c_1\) large enough that every constructed
coordinate is nonnegative, and set
\begin{equation}
\begin{aligned}
  r(c_{j+1}) &= r(c_j)+L,&
  q(c_{j+1}) &= q(c_j)+L+s_jd_1,\\
  r(x_{j,t}) &= r(c_j)+T+t\tau,&
  q(x_{j,t}) &= q(c_j)+T-s_jd_2+t\tau.
\end{aligned}
\label{eq:chaining-quadratic-coordinate-witness}
\end{equation}
This places every branch before \(c_{j+1}\) in both coordinates and gives the
diagonals in Equation~\ref{eq:chaining-quadratic-diagonals}. Both
\(c_j\)-to-\(c_{j+1}\) advances are below \(X\). For every anchor in group
\(j-2\) or earlier, the reference and query advances to group \(j\) exceed
\(X\). Two consecutive orange branches differ by \(d_1+2d_2>B\). For the
numerical example above, one valid coordinate choice is \(h_0=1000\),
\(T=338\), \(\tau=88\), and \(L=4000\).

Choose scores \(F_X,F_C,F_E,F_G\) with
\begin{equation}
\begin{gathered}
  k<F_X=F_C+\Lambda(d_2)<F_C<S,\\
  S\leq F_E=F_X+\ell k,\qquad
  \ell\geq K,\qquad \ell k\geq S,\\
  \ell\leq\texttt{max\_skip},\qquad
  \ell+2\leq\texttt{max\_iter},\\
  F_G-F_X>D.
\end{gathered}
\label{eq:chaining-quadratic-score-conditions}
\end{equation}
A prefix of 35 anchors on one diagonal has scores
\(15,30,\ldots,525\), so set \(F_G=525\). Follow it with 27 links whose
diagonal difference is \(d_2=250\). Each link contributes \(-18\), so these
scores decrease to \(39=F_C\). Give every connecting-path link reference and
query advances greater than \(X/2\) and at most \(X\). An anchor two links back
is then outside the distance limit. Every connecting-path score is at least
\(F_C\), and
\[
  F_G-F_X=525-21=504>D=500.
\]

The first score-reducing anchor can inspect several anchors in the high-scoring
prefix. The final prefix anchor \(h^\star\) has the largest prefix score, and
all of these candidates have the same diagonal penalty. Therefore
\(h^\star\) is its predecessor. For later score-reducing links, the
greater-than-\(X/2\) advances exclude every candidate except the immediate
predecessor.

The prefix end \(h^\star\) is the first backtracking candidate and marks the
high-scoring prefix. A score-reducing anchor with score at least \(S\) can also
enter the candidate list before an anchor \(e_j\). All of its predecessors on
that path have larger scores. Its cut search therefore returns that
score-reducing anchor, and
the permanent loop marks no score-reducing anchor. The path remains available
to every later branch search.

The forward pass makes the choices shown in
Figure~\ref{fig:chaining-quadratic-anchor-family}. For \(x_{j,0}\), candidate
\(c_j\) gives
\[
  F_C+\Lambda(d_2)=F_X>k.
\]
If \(c_{j-1}\) is inside the distance limit, it gives
\[
  F_C+\Lambda(d_1+d_2)<F_X.
\]
Every anchor in the preceding orange branch is outside the bandwidth because
its diagonal difference is \(d_1+2d_2>B\). Anchors from older groups are
outside the distance limit. Thus \(p(x_{j,0})=c_j\).

For \(t>0\), the immediate orange predecessor gives \(F_X+tk\). Every earlier
anchor in that branch has a lower score. Candidates outside the branch are
bounded as above, so each later branch anchor chooses its immediate
predecessor.

For \(c_{j+1}\), candidate \(c_j\) gives \(F_C\). The highest-scoring candidate
in the current orange branch is \(e_j\), which gives
\[
  F_E+\Lambda(d_1+d_2)<F_C.
\]
Every candidate in group \(j-1\) or earlier is outside the distance limit. The
scan reaches \(c_j\)
because \(\ell\leq\texttt{max\_skip}\) and
\(\ell+2\leq\texttt{max\_iter}\). Therefore \(p(c_{j+1})=c_j\), and every
backbone score remains \(F_C\).

\begin{figure}[H]
\centering
\begin{tikzpicture}[
  x=1cm,
  y=1cm,
  font=\footnotesize,
  proof anchor/.style={
    circle,
    draw=black!70,
    fill=white,
    minimum size=.46cm,
    inner sep=0pt
  },
  proof backbone/.style={
    proof anchor,
    draw=chainBlue!90!black,
    fill=chainBlue!16
  },
  proof branch/.style={
    proof anchor,
    draw=chainOrange!95!black,
    fill=chainOrange!18
  },
  proof prefix/.style={
    proof anchor,
    draw=chainGreen!80!black,
    fill=chainGreen!15
  },
  predecessor/.style={-{Latex[length=1.7mm]},thick,draw=black!68},
  proof note/.style={
    draw=chainGray!75!black,
    fill=chainGray!8,
    rounded corners=2pt,
    align=center,
    inner sep=3pt
  }
]
  \node[anchor=west,font=\bfseries] at (.35,5.35)
    {(a) Predecessors produced by forward chaining};

  \node[proof prefix] (hp) at (.85,3.20) {$h^\star$};
  \node[proof anchor] (h2) at (1.62,3.20) {};
  \node[proof anchor] (h1) at (2.34,3.20) {};
  \node[proof backbone] (c1) at (3.25,3.20) {$c_1$};
  \node[proof backbone] (c2) at (4.65,3.20) {$c_2$};
  \node[text=black!55] at (5.62,3.20) {$\cdots$};
  \node[proof backbone] (cj) at (6.65,3.20) {$c_j$};
  \node[text=black!55] at (7.72,3.20) {$\cdots$};
  \node[proof backbone] (cm) at (8.78,3.20) {$c_m$};

  \draw[predecessor] (h1)--(h2);
  \draw[predecessor] (h2)--(hp);
  \draw[predecessor] (c1)--(h1);
  \draw[predecessor,draw=chainBlue!85!black] (c2)--(c1);
  \draw[predecessor,draw=chainBlue!85!black] (cj)--(c2);
  \draw[predecessor,draw=chainBlue!85!black] (cm)--(cj);

  \foreach \base/\idx in {3.25/1,4.65/2,6.65/j,8.78/m} {
    \node[proof branch] (x\idx) at (\base,4.02) {$x_{\idx,0}$};
    \node[proof branch] (e\idx) at (\base,4.82) {$e_{\idx}$};
    \draw[predecessor,dashed,draw=chainOrange!90!black] (e\idx)--(x\idx);
    \node[fill=white,inner sep=.5pt] at (\base,4.42) {$\vdots$};
  }
  \draw[predecessor,draw=chainOrange!90!black] (x1)--(c1);
  \draw[predecessor,draw=chainOrange!90!black] (x2)--(c2);
  \draw[predecessor,draw=chainOrange!90!black] (xj)--(cj);
  \draw[predecessor,draw=chainOrange!90!black] (xm)--(cm);

  \node[anchor=south,text=chainGreen!60!black] at (hp.north) {$F_G$};
  \node[anchor=north,text=chainBlue!80!black] at (cj.south) {$F_C$};
  \node[anchor=west,text=chainOrange!90!black] at (xj.east) {$F_X$};
  \node[anchor=west,text=chainOrange!90!black] at (ej.east) {$F_E$};

  \draw[decorate,decoration={brace,amplitude=4pt},chainOrange!90!black]
    (9.25,4.00)--(9.25,4.84)
    node[midway,right=4pt,text=black!75]
      {$\ell$ marked anchors};
  \node[proof note] at (11.45,3.48)
    {$\operatorname{ChooseCut}(e_j)=x_{j,0}$\\
     $c_j,\ldots,c_1$ are cleared};

  \draw[black!18] (.35,2.45)--(12.85,2.45);
  \node[anchor=west,font=\bfseries] at (.35,2.08)
    {(b) Predecessor links revisited by searches from \(e_j\)};

  \node[anchor=east,text=black!65] at (2.12,1.48) {$e_1$};
  \node[anchor=east,text=black!65] at (2.12,.82) {$e_j$};
  \node[anchor=east,text=black!65] at (2.12,.16) {$e_m$};
  \node[font=\scriptsize,text=chainOrange!90!black] at (2.80,1.78)
    {$\ell$ branch links};
  \foreach \x in {2.55,3.05}
    \fill[chainOrange!90!black] (\x,1.48) circle[radius=2.1pt];
  \fill[chainBlue!85!black] (3.55,1.48) circle[radius=2.1pt];

  \foreach \x in {2.55,3.05}
    \fill[chainOrange!90!black] (\x,.82) circle[radius=2.1pt];
  \foreach \x in {3.55,4.05,4.55,5.05,5.55}
    \fill[chainBlue!85!black] (\x,.82) circle[radius=2.1pt];
  \node[text=black!55] at (6.05,.82) {$\cdots$};
  \fill[chainBlue!85!black] (6.55,.82) circle[radius=2.1pt];

  \foreach \x in {2.55,3.05}
    \fill[chainOrange!90!black] (\x,.16) circle[radius=2.1pt];
  \foreach \x in {3.55,4.05,4.55,5.05,5.55,6.05,6.55,7.05,7.55}
    \fill[chainBlue!85!black] (\x,.16) circle[radius=2.1pt];
  \node[text=black!55] at (8.05,.16) {$\cdots$};
  \fill[chainBlue!85!black] (8.55,.16) circle[radius=2.1pt];

  \node[proof note,text width=3.55cm] at (10.82,.82) {
    orange: links within one branch\\
    blue: backbone links visited again
  };
\end{tikzpicture}
\caption{Predecessor structure used in the lower bound. Arrows point from
each anchor to its predecessor. Panel (a) shows the high-scoring prefix anchor
\(h^\star\), backbone anchors \(c_j\), and branch anchors
\(x_{j,0},\ldots,e_j\). A dashed orange segment abbreviates \(\ell\)
predecessor links. Backtracking marks the \(\ell\) branch anchors after
\(x_{j,0}\) and leaves \(x_{j,0}\) unmarked. Panel (b) shows that the search
from \(e_j\) revisits the first \(j\) backbone anchors.}
\label{fig:chaining-quadratic-anchor-family}
\end{figure}

For a search starting at \(e_j\), the predecessor walk is
\[
  e_j\to\cdots\to x_{j,0}\to c_j\to c_{j-1}\to\cdots\to c_1
  \to\text{the fixed prefix path}.
\]
At \(x_{j,0}\), \textsc{ChooseCut} records the decrease
\(F_E-F_X=\ell k\). Every backbone and connecting-path score is at least
\(F_C>F_X\). No predecessor visited after \(x_{j,0}\) changes the recorded
cut. When the walk reaches score
\(F_G\), the decrease from the recorded maximum is
\[
  (F_E-F_X)-(F_E-F_G)=F_G-F_X>D.
\]
The walk stops and returns \(x_{j,0}\). The permanent pass marks exactly the
\(\ell\) branch anchors after this cut. The temporary-mark cleanup inside
\textsc{ChooseCut} clears the shared backbone, which allows every later branch
search to visit it again. Anchors \(e_j\) with equal scores may be processed
in either order, and \textsc{ChooseCut} still returns \(x_{j,0}\). If the
high-scoring prefix has already been marked,
the walk stops at that mark. The marked prefix score remains above \(F_X\), so
the returned cut is unchanged.

The candidate list contains every anchor whose score is at least \(S\). Some
internal branch anchors can enter this list. Each has a lower score than
\(e_j\). Processing \(e_j\) first marks every such internal candidate in its
branch, so it is skipped later. The output from \(e_j\) contains exactly
\(\ell\) anchors and has extracted score
\[
  F_E-F_X=\ell k.
\]
It is accepted because \(\ell\geq K\) and \(\ell k\geq S\). Branch \(j\)
follows \(j+O(1)\) predecessor links.
Consequently,
\begin{equation}
  \sum_{j=1}^{m}\bigl(j+O(1)\bigr)=\Theta(m^2),
  \qquad
  n=(\ell+2)m+O(1)=\Theta(m),
  \label{eq:chaining-quadratic-retained-anchor-bound}
\end{equation}
and scalar backtracking takes \(\Theta(n^2)\) time in the worst case. Both
scalar backtracking passes perform the same asymptotic number of predecessor
link traversals, so the second pass changes only the constant factor.

Equation~\ref{eq:chaining-conserved-block-existence} proves that enough fixed
blocks exist for finite examples with at most 185 conserved blocks.
Equation~\ref{eq:chaining-periodic-padding-existence} supplies long padding
with selected minimizers and no anchors.
Equation~\ref{eq:chaining-complete-filler-event} gives a probability bound for
the remaining random fillers. Together, these equations establish the finite
sequence template in Figure~\ref{fig:chaining-quadratic-sequence-template}.

The asymptotic lower bound is deterministic. It depends on the retained
anchors and the \(f\) and \(p\) arrays consumed by backtracking. This section
does not claim an unbounded fixed-\(k\) FASTA family. Such a family would need
more distinct retained \(k\)-mer values than the finite \(k\)-mer alphabet eventually
provides. Allowing \(k\) to grow would also change the score function. The
\(\Theta(n^2)\) result is therefore a worst-case theorem for the retained
anchor, score, and predecessor arrays. Positive minimizer density does not
prevent finite retained-anchor arrays that require quadratic backtracking.

\coderef{minimap2/lchain.c}{Minimap2 chaining link score}{comput_sc}{113}
\coderef{minimap2/lchain.c}{Minimap2 cut search}{mg_chain_bk_end}{9}
\coderef{minimap2/lchain.c}{Minimap2 scalar backtracking}{mg_chain_backtrack}{27}